\documentclass[10pt,a4paper]{amsart}
\usepackage{amsmath,amssymb,amsthm,graphicx}
\usepackage{verbatim}
\newenvironment{reference}{\comment}{\endcomment}
%\newenvironment{reference}{}{}
\newenvironment{slogan}{\comment}{\endcomment}
\newenvironment{history}{\comment}{\endcomment}

% For commutative diagrams we use Xy-pic
\usepackage[all]{xy}

% We use 2cell for 2-commutative diagrams.
\xyoption{2cell}
% Theorem environments.
%
\theoremstyle{plain}
\newtheorem{theorem}[subsection]{Theorem}
\newtheorem{proposition}[subsection]{Proposition}
\newtheorem{lemma}[subsection]{Lemma}

\theoremstyle{definition}
\newtheorem{definition}[subsection]{Definition}
\newtheorem{example}[subsection]{Example}
\newtheorem{exercise}[subsection]{Exercise}
\newtheorem{situation}[subsection]{Situation}

\theoremstyle{remark}
\newtheorem{remark}[subsection]{Remark}
\newtheorem{remarks}[subsection]{Remarks}

\numberwithin{equation}{subsection}

% Macros
%
\def\lim{\mathop{\mathrm{lim}}\nolimits}
\def\colim{\mathop{\mathrm{colim}}\nolimits}
\def\Spec{\mathop{\mathrm{Spec}}}
\def\Hom{\mathop{\mathrm{Hom}}\nolimits}
\def\Ext{\mathop{\mathrm{Ext}}\nolimits}
\def\SheafHom{\mathop{\mathcal{H}\!\mathit{om}}\nolimits}
\def\SheafExt{\mathop{\mathcal{E}\!\mathit{xt}}\nolimits}
\def\Sch{\mathit{Sch}}
\def\Mor{\mathop{\mathrm{Mor}}\nolimits}
\def\Ob{\mathop{\mathrm{Ob}}\nolimits}
\def\Sh{\mathop{\mathit{Sh}}\nolimits}
\def\NL{\mathop{N\!L}\nolimits}
\def\CH{\mathop{\mathrm{CH}}\nolimits}
\def\proetale{{pro\text{-}\acute{e}tale}}
\def\etale{{\acute{e}tale}}
\def\QCoh{\mathit{QCoh}}
\def\Ker{\mathop{\mathrm{Ker}}}
\def\Im{\mathop{\mathrm{Im}}}
\def\Coker{\mathop{\mathrm{Coker}}}
\def\Coim{\mathop{\mathrm{Coim}}}

% Boxtimes
%
\DeclareMathSymbol{\boxtimes}{\mathbin}{AMSa}{"02}

%
% Macros for moduli stacks/spaces
%
\def\QCohstack{\mathcal{QC}\!\mathit{oh}}
\def\Cohstack{\mathcal{C}\!\mathit{oh}}
\def\Spacesstack{\mathcal{S}\!\mathit{paces}}
\def\Quotfunctor{\mathrm{Quot}}
\def\Hilbfunctor{\mathrm{Hilb}}
\def\Curvesstack{\mathcal{C}\!\mathit{urves}}
\def\Polarizedstack{\mathcal{P}\!\mathit{olarized}}
\def\Complexesstack{\mathcal{C}\!\mathit{omplexes}}
% \Pic is the operator that assigns to X its picard group, usage \Pic(X)
% \Picardstack_{X/B} denotes the Picard stack of X over B
% \Picardfunctor_{X/B} denotes the Picard functor of X over B
\def\Pic{\mathop{\mathrm{Pic}}\nolimits}
\def\Picardstack{\mathcal{P}\!\mathit{ic}}
\def\Picardfunctor{\mathrm{Pic}}
\def\Deformationcategory{\mathcal{D}\!\mathit{ef}}
\usepackage{cleveref}
\crefname{originallemma}{original lemma}{original lemmas}
\Crefname{originallemma}{Original lemma}{Original lemmas}
\makeatletter
\newlabel{algebra-lemma-finitely-generated-pure-ideal}{{05KK}{source}{}{sourcepointer-1}{}}
\newlabel{algebra-lemma-finitely-generated-pure-ideal@cref}{{[originallemma][0][]05KK}{[1][0][]source}{}{}{}}
\makeatother
\crefname{originallemma}{original lemma}{original lemmas}
\Crefname{originallemma}{Original lemma}{Original lemmas}
\makeatletter
\newlabel{algebra-lemma-pure}{{04PS}{source}{}{sourcepointer-2}{}}
\newlabel{algebra-lemma-pure@cref}{{[originallemma][0][]04PS}{[1][0][]source}{}{}{}}
\makeatother
\crefname{originalsection}{original section}{original sections}
\Crefname{originalsection}{Original section}{Original sections}
\makeatletter
\newlabel{section-finite-flat-not-projective}{{052H}{source}{}{sourcepointer-3}{}}
\newlabel{section-finite-flat-not-projective@cref}{{[originalsection][0][]052H}{[1][0][]source}{}{}{}}
\makeatother
\title{Flat finite type ring maps and proper flat morphisms with controlled fibres that are not of finite presentation}
\author{The Stacks Project Authors}
\date{stacks-project:examples.tex:local-snapshot-20261012; GNU FDL 1.2 or later}
\begin{document}
\maketitle
\noindent\textit{Editorial note.} All delivered mathematical text is the author's own, copied from the pinned Stacks Project chapter file examples.tex; no statement or proof step was rewritten, reordered or removed.\par\medskip
\noindent\textit{Editorial note.} The counted proof is the author's complete argument: The four context paragraphs map one-to-one onto the four enumerated items; the (R, I) base data is the delivered L48315 example, making the chapter self-contained.\par\medskip
\noindent\textit{Editorial note.} The label lemma-finite-type-flat-not-finitely-presented names the counted statement, and the proof environment immediately following it in the source is delivered verbatim as its proof.\par\medskip
\noindent\textit{Editorial note.} The counted statement is carried by the author's own lemma environment in the chapter source.\par\medskip
\noindent\textit{Allowed standard prerequisites (editorial).}\begin{itemize}
\item Context span 3050-3059: R is a ring with an ideal I such that R/I is a finite flat module but not projective (explicit example delivered as card L48315, section-finite-flat-not-projective); this means I is not finitely generated (Algebra, finitely-generated-pure-ideal) and I = I\textasciicircum{}2 (Algebra, pure); the base scheme is S = Spec(R)
\item Context span 3062-3065: the ring map R -> R + (R/I)epsilon with epsilon\textasciicircum{}2 = 0 is finite, flat, not of finite presentation, with all fibre rings complete intersections and geometrically irreducible
\item Context span 3068-3073: A = R[x, y]/(xy, ay for a in I) is as an R-module the direct sum of R y\textasciicircum{}i (i >= 0) and (R/I) x\textasciicircum{}j (j > 0); R -> A is flat finite type not of finite presentation, each fibre ring isomorphic to kappa(p)[x, y]/(xy) or kappa(p)[x]
\item Context span 3076-3082: B = R[X\_0, X\_1, X\_2]/(X\_1 X\_2, a X\_2 for a in I) gives Proj(B) -> S proper flat not of finite presentation with fibres P\textasciicircum{}1\_s or the vanishing locus of X\_1 X\_2 in P\textasciicircum{}2\_s (a projective nodal curve of arithmetic genus 0)
\item Context span 3085-3090: M = R + R + R/I, B = Sym\_R(M), X = Proj(B): X -> S is proper flat not of finite presentation with geometric fibres P\textasciicircum{}1\_s or P\textasciicircum{}2\_s, in particular smooth and geometrically irreducible
\end{itemize}\medskip
\ifdefined\hypertarget\hypertarget{sourcepointer-1}{}\fi
\noindent\textit{Source reference.} Original lemma 05KK: sources/stacks-chow/tags.txt, source lines 7204--7204. Prerequisite pointer invoked at 3118; it names a standard fact of a sibling chapter that the counted argument uses, and every step of the counted proof built on it is retained verbatim in the delivered body.\par\medskip
\ifdefined\hypertarget\hypertarget{sourcepointer-2}{}\fi
\noindent\textit{Source reference.} Original lemma 04PS: sources/stacks-chow/tags.txt, source lines 6125--6125. Prerequisite pointer invoked at 3120; it names a standard fact of a sibling chapter that the counted argument uses, and every step of the counted proof built on it is retained verbatim in the delivered body.\par\medskip
\ifdefined\hypertarget\hypertarget{sourcepointer-3}{}\fi
\noindent\textit{Source reference.} Original section 052H: sources/stacks-examples/pinned-original-source.tex, source lines 2620--2620. Prerequisite pointer invoked at 3115; it names a standard fact of the same chapter that the counted argument uses, and every step of the counted proof built on it is retained verbatim in the delivered body.\par\medskip
Let $R$ be a ring which has an ideal $I$ such that $R/I$ is a finite
flat module but not projective, see
Section \ref{section-finite-flat-not-projective}
for an explicit example. Note that this means that $I$ is not
finitely generated, see
Algebra, Lemma \ref{algebra-lemma-finitely-generated-pure-ideal}.
Note that $I = I^2$, see
Algebra, Lemma \ref{algebra-lemma-pure}.
The base ring in our examples will be
$R$ and correspondingly the base scheme $S = \Spec(R)$.

\medskip\noindent
Consider the ring map $R \to R \oplus R/I\epsilon$ where
$\epsilon^2 = 0$ by convention. This is a finite, flat
ring map which is not of finite presentation. All the fibre rings
are complete intersections and geometrically irreducible.

\medskip\noindent
Let $A = R[x, y]/(xy, ay; a \in I)$. Note that as an $R$-module
we have $A = \bigoplus_{i \geq 0} Ry^i \oplus \bigoplus_{j > 0} R/Ix^j$.
Hence $R \to A$ is a flat finite type ring map which is not of finite
presentation. Each fibre ring
is isomorphic to either $\kappa(\mathfrak p)[x, y]/(xy)$
or $\kappa(\mathfrak p)[x]$.

\medskip\noindent
We can turn the previous example into a projective morphism by
taking $B = R[X_0, X_1, X_2]/(X_1X_2, aX_2; a \in I)$.
In this case $X = \text{Proj}(B) \to S$ is a proper flat morphism
which is not of finite presentation such that for each $s \in S$
the fibre $X_s$ is isomorphic either to $\mathbf{P}^1_s$ or to
the closed subscheme of $\mathbf{P}^2_s$ defined by the vanishing of
$X_1X_2$ (this is a projective nodal curve of arithmetic genus $0$).

\medskip\noindent
Let $M = R \oplus R \oplus R/I$. Set $B = \text{Sym}_R(M)$ the
symmetric algebra on $M$. Set $X = \text{Proj}(B)$.
Then $X \to S$ is a proper flat morphism, not of finite presentation
such that for $s \in S$ the geometric fibre is isomorphic to either
$\mathbf{P}^1_s$ or $\mathbf{P}^2_s$. In particular these fibres are
smooth and geometrically irreducible.

\begin{lemma}
\label{lemma-finite-type-flat-not-finitely-presented}
There exist examples of
\begin{enumerate}
\item a flat finite type ring map with geometrically irreducible
complete intersection fibre rings which is not of finite presentation,
\item a flat finite type ring map with geometrically connected,
geometrically reduced, dimension 1, complete intersection fibre rings
which is not of finite presentation,
\item a proper flat morphism of schemes $X \to S$ each of whose fibres
is isomorphic to either $\mathbf{P}^1_s$ or to the vanishing locus of
$X_1X_2$ in $\mathbf{P}^2_s$ which is not of finite presentation, and
\item a proper flat morphism of schemes $X \to S$ each of whose
fibres is isomorphic to either $\mathbf{P}^1_s$ or $\mathbf{P}^2_s$
which is not of finite presentation.
\end{enumerate}
\end{lemma}

\begin{proof}
See discussion above.
\end{proof}

\end{document}
